% Section 1: The Ideal Switch with Realistic Components

\subsection{The DPT Circuit}

We begin with a \textbf{Double Pulse Test (DPT)} circuit—one of the simplest configurations used to characterize semiconductor switches. The circuit consists of:

\begin{itemize}
    \item A DC voltage source $\vdc = 800\,\si{\volt}$
    \item A load inductor $L$ with series resistance $R_L$ (parasitic winding resistance)
    \item A switch $Q$ (our MOSFET)
    \item A freewheeling diode $D$ (body diode of the switch, modeled as piecewise linear)
    \item Gate signal $\vgs(t)$ that controls the switch
\end{itemize}

\noindent The circuit diagram is shown in Figure~\ref{fig:dpt_circuit}.

\begin{figure}[h]
    \centering
    \begin{tikzpicture}[
        scale=1.2,
        circuit logic US,
        thick
    ]
        % Supply
        \node[draw, circle, minimum size=0.8cm] (vsrc) at (0,4) {$+$};
        \node at (-0.15, 3.7) {$V_{DC}$};

        % Load inductor
        \node[draw] (L) at (1.5, 4) [rectangle, minimum width=0.8cm, minimum height=0.5cm] {L};

        % Current arrow
        \node[below=0.3cm of L] (Ilabel) {$I_L$};

        % Diode
        \node[draw] (D) at (3.5, 4) [rectangle, minimum width=0.8cm, minimum height=0.5cm] {D};

        % Switch Q
        \node[draw] (Q) at (3.5, 2) [rectangle, minimum width=0.8cm, minimum height=0.5cm] {Q};

        % Voltages
        \draw[->] (3.8, 3.3) -- (3.8, 2.8) node[right, font=\small] {$v_{ds}$};
        \draw[->] (3.2, 1.5) -- (3.2, 0.8) node[left, font=\small] {$i_d$};
        \draw[-] (0, 0.2) -- (0, 0) node[ground] {};

        % Gate signal
        \draw (-1, 0.5) -- (1.5, 0.5) node[above, font=\small] {$v_{gs}(t)$};

        % Connections
        \draw (vsrc) -- (L);
        \draw (L) -- (D);
        \draw (D) -- (Q);
        \draw (Q) -- (0, 2);
        \draw (0, 2) -- (0, 0);
        \draw (3.5, 4) -- (3.5, 2.8);
        \draw (3.5, 1.2) -- (3.5, 0);
        \draw (3.5, 0) -- (0, 0);

    \end{tikzpicture}
    \caption{Double Pulse Test (DPT) circuit: DC source + load inductor + MOSFET + body diode.}
    \label{fig:dpt_circuit}
\end{figure}

\subsection{The Piecewise Linear MOSFET Model}

In this first section, we use a \textbf{piecewise linear model} for the switch $Q$ with realistic turn-on/off behavior:

\subsubsection{MOSFET Turn-On/Off Logic}

\begin{enumerate}
    \item \textbf{ON state:} When $\vgs \geq \vth$ (gate voltage reaches or exceeds threshold voltage):
    \[
        \ids = \frac{\vds}{\ron}, \quad \text{where } \ron \approx 0
    \]
    The switch conducts with near-zero ``on-resistance''. For simplicity, we approximate $\vds \approx 0$ when conducting.

    \item \textbf{OFF state:} When $\vgs < \vth$, the switch is reverse-biased:
    \[
        \ids = 0
    \]
    No current flows through $Q$. The full supply voltage appears across $\vds$ (minus the diode voltage).

    \item \textbf{No capacitance:} In this first section, the switch has no parasitic energy storage (no $C_{DS}$, no $C_{GS}$, no $C_{GD}$). These will be added later.

    \item \textbf{Threshold voltage:} The switch turns ON when $\vgs$ crosses the threshold:
    \[
        \vth \approx 2\text{--}3\,\si{\volt} \quad \text{(typical SiC MOSFET)}
    \]
\end{enumerate}

\subsubsection{Body Diode: Piecewise Linear Model (PWL)}

The freewheeling diode $D$ is modeled using the \textbf{piecewise linear (PWL)} model with three components:
\begin{itemize}
    \item An ideal diode (conducts only in forward direction)
    \item A built-in voltage source $V_{bo}$ (forward voltage drop threshold)
    \item A dynamic resistance $r_d$ (diode on-resistance)
\end{itemize}

The PWL diode equations are:

\begin{enumerate}
    \item \textbf{Forward biased} ($i_d > 0$, diode conducts):
    \[
        \vds = V_{bo} + i_d \, r_d
    \]
    where $V_{bo}$ is the built-in voltage (e.g., $0.6\text{--}1.2\,\si{\volt}$ for SiC) and $r_d$ is the forward resistance.

    \item \textbf{Reverse biased} ($i_d = 0$, diode blocks):
    \[
        i_d = 0, \quad \vds < V_{bo}
    \]
    No current flows. The voltage $\vds$ is determined by the circuit (can be negative).
\end{enumerate}

This PWL model replaces the more complex Shockley exponential equation with a simple piecewise linear relationship that is easy to solve analytically.

\noindent\textbf{Key insight:} Once we replace the junction diode with this PWL equivalent, the circuit becomes purely resistive and algebraic—no implicit equations needed.

\subsection{Circuit Equations with Ideal Switch}

With the ideal switch, the circuit behavior is \textbf{purely algebraic} (no differential equations yet):

\begin{enumerate}
    \item \textbf{When Q is ON} ($\vgs > \vth$):
    \[
        \vds = 0, \quad \text{inductor voltage is } \vdc
    \]
    The inductor current $I_L$ increases linearly (we'll address this in the next section).

    \item \textbf{When Q is OFF} ($\vgs < \vth$):
    \[
        \ids = 0, \quad \text{current commutates to diode } D
    \]
    The diode conducts, and $\vds = -V_D$ (the negative diode forward drop, which we assume is zero).

    \item \textbf{Body diode forward drop:}
    \[
        \text{If diode conducts:} \quad \vds = 0 \quad (\text{ideal diode})
    \]
\end{enumerate}

\subsection{The Problem with This Model}

The ideal switch model predicts \textbf{instantaneous} transitions:
\begin{itemize}
    \item Switch turns ON → $\vds$ drops to $0$ instantly
    \item Switch turns OFF → $\vds$ jumps to $\vdc$ instantly
    \item No intermediate voltages, no oscillations, no overshoot
\end{itemize}

But in real hardware, we observe:
\begin{itemize}
    \item \textbf{Switching transients:} $\vds$ takes time to rise or fall (not instantaneous)
    \item \textbf{Voltage overshoot:} $\vds$ can exceed $\vdc$ during turn-off
    \item \textbf{Ringing and oscillations:} High-frequency oscillations appear in $\vds$ and $\ids$
\end{itemize}

These phenomena are \textbf{not captured} by an ideal switch model with zero/infinite resistance.

\subsection{What's Missing?}

\boxed{\text{Energy storage}}

The ideal switch model assumes resistive transitions. But real switching transients involve \textbf{energy flows between electric and magnetic fields}. The inductor $L$ stores magnetic energy, and parasitic capacitances store electric energy.

\vspace{1em}
\noindent\textbf{Key question for the next section:}
\begin{center}
    \fbox{\parbox{0.8\textwidth}{
        If the inductor current cannot change instantly (because $L \frac{di_L}{dt} = V$),
        what happens to the voltage $\vds$ when the switch turns OFF?
        How do we model the resulting voltage rise?
    }}
\end{center}

\vspace{1em}
\noindent\textit{In the next section, we introduce the load inductor as an explicit state variable and solve the resulting differential equation.}
